\section{Three exact arithmetic admissibility checks}
\label{sec:arithmetic-admissibility}

The preceding coefficient constructions are actual constructions.
They must still be distinguished from the additional hypotheses of
a complete step-10 divisibility application. The following three
checks prove incompatibilities of specific supplied specializations;
they are stronger than unsuccessful searches for an input file.

\begin{proposition}[SF-Q6-01: the numerical examples are not global records]
There are no global datum $D:\texttt{D7}$, size datum
$L:\texttt{SizeData}_2(D)$, and nonnegative integers
$\Delta,h_0$ for which \texttt{ofRecord D L Delta h0} has
the rank and host fields of \texttt{numEx} or \texttt{numBig}.
\end{proposition}
\begin{proof}
The definitions \texttt{hostEx}, \texttt{numEx} and
\texttt{ofRecord}, with the equalities \texttt{hostEx\_q} and
\texttt{hostEx\_d}, occur in
\path{Formalism/LocalToGlobal/Divisibility/NumbersLemmas.lean}.
The definition \texttt{numBig} in the adjacent
\path{NumbersLink.lean} copies these rank and host fields:
\[
 (k,m,t,r,f)=(8,2,15,9,8),\qquad(q,d)=(11,7).
\]
An actual global datum has \texttt{D.host : Host D.R}.
The \texttt{qd\_big} field of \texttt{Host}, in
\path{Formalism/LocalRingCardinality/DataHost.lean}, requires
$4tr+1<q$ and $4tr+1<d$. These become $541<11$ and
$541<7$, a contradiction. The definition \texttt{ofRecord}
copies these global fields, so none of the other choices repairs
them. This does not invalidate the numerical inequalities proved
by \texttt{numFactsBig} on its declared numerical domain.
\end{proof}

The supplied \texttt{Examples}\(_2\)\texttt{.H0} in
\path{Formalism/LocalRingCardinality/SpanPowerMembership/ExampleLayerExistsGenerators.lean}
has the permissible host values $(q,d)=(571,599)$. That is a full
host, not a complete global row-and-period datum. In particular,
the later level-prime congruences of \texttt{Pre} are additional
conditions; the host alone is not a substitute for them.

\begin{proposition}[SF-Q6-02: the inhabited small-weight jets fail the size condition]
Suppose a global datum has $k=8$, a nonempty top set
\texttt{host.TT}, and last Harish--Chandra coordinate
$u_7=-1/2$ at every top embedding. Then it admits no
$\texttt{SizeData}_2$. In particular this applies to the
specific constructed witness of IF-11.
\end{proposition}
\begin{proof}
Choose $t\in\texttt{host.TT}$ and the corresponding top embedding,
using \texttt{D.eTau} and \texttt{tau1}. The field
\texttt{SizeData}\(_2\)\texttt{.top\_u}, in
\path{Formalism/LocalRingCardinality/} followed by
\texttt{DataParams}\(_2\)\texttt{.lean}, lines 10--101,
would give
\[
 -\tfrac12=-b_T(t)-\tfrac72,\qquad b_T(t)=-3.
\]
The \texttt{bT\_ge} field instead gives $b_T(t)\ge b'>0$.
Here the scale definitions are
$b=\exp(x^\varepsilon)$, $a=xb$ and
$b'=a/\exp(x^{\varepsilon/3})$, with $x>1$, so the strict
positivity is immediate. This is a contradiction. In the selected
IF-11 construction the top set is $\{0\}$ and the endpoint is
$-1/2$. The claim concerns that traced witness; the unrefined
existential theorem does not require every possible witness to
have those parameters.
\end{proof}

\begin{proposition}[SF-Q6-03: the asymptotic rank window does not repair these ranks]
For $(k,m,t,r,f)=(8,2,15,9,8)$, no scale satisfies
\texttt{GlobalTheory.EX.RankVal}. More generally,
$m=2$ together with \texttt{RankVal} implies $x<81$.
\end{proposition}
\begin{proof}
The last condition of \texttt{RankVal}, in
\path{Development/AbelianVarietyTheory/Definitions/GlobalTwisting/PairRowIffQdRows.lean},
lines 17--32, is the equality of nonnegative integers
\[
 r+t=8fr+4+\texttt{TOv}(s)\,\texttt{nWv}(R,s).
\]
The fixed ranks give $24\ge580$, which is impossible.
For the second assertion, the same definition requires
\[
 \sqrt{x\delta_1}\le m,\qquad
 \delta_1=x^{-\varepsilon/8},\qquad
 x>1,\quad 0<\varepsilon<1/52.
\]
Since $1-\varepsilon/8>1/2$,
$x\delta_1=x^{1-\varepsilon/8}\ge\sqrt{x}$, hence
$m\ge x^{1/4}$. If $x\ge81$, then $m\ge3$,
contradicting $m=2$. The supplied general implication is
\texttt{NodeGap.le\_m}, in
\path{Development/AbelianVarietyTheory/Proofs/IwasawaDecomposition/EventuallyNodeGapGeRankPow.lean},
lines 20--35. The short argument here records its specified
two-step consequence and the separate fixed-rank contradiction.
\end{proof}

The constructor \texttt{exists\_sizeData\_explicit} in
\path{Proofs/GlobalInterpolation/ExistsExplicitParameters.lean}
(relative to \path{Development/AbelianVarietyTheory/}) assumes
\texttt{RankVal}; the assembled route in
\path{Proofs/IwasawaDecomposition/SizeParametersNonempty.lean}
uses \texttt{Ctx.ok.rank}. It therefore cannot supply the
missing size datum for IF-11's fixed ranks. An unbounded-scale
family with that rank window cannot keep $m=2$ either.
The abstract types \texttt{SizeData}\(_2\) and
\texttt{DivInputs} do not impose \texttt{RankVal} universally.
A direct construction with a different host or weights and all
of its bounds verified remains a distinct approach.

\paragraph{Arithmetic and geometric meaning.}
These restrictions precede the intended arithmetic application.
The numerical polynomial representations still define exterior and
$\lambda$-functors over $\operatorname{Spec}\mathbb Z$, and the
IF-11 family still supplies actual CM-normalized algebraic jets.
The excluded host and size values do not enter the abelian
$C+I$ theorem or the corrected geometric height theorem. They
are additional requirements of the intended arithmetic lattice
construction, with its local denominators. Neither a rational
$K_0$ character, an $\ell$-adic weight, nor a motivic local
$L$-factor supplies these missing numerical equalities or
inequalities. Their relationship to cote 16 scans 014--018 is
therefore through the explicitly specified coefficient and
operator models, not through an identification of those
arithmetic constraints with a Lefschetz conjecture.
